% !PRINTED 440
\weberpage{440}
%<<CONT
so erhalten wir als Bedingung dafür, dass der Punkt $y$ mit dem Punkte $x$ zusammenfällt, da ein Punkt durch die \gesperrt{Verhältnisse} seiner Coordinaten eindeutig bestimmt ist, die, dass für einen passend bestimmten Co\"efficienten $\lambda$:
\weq{3}{\begin{aligned}
\lambda x_1 &= \alpha\,x_1 + \beta\,x_2 + \gamma\,x_3 \\
\lambda x_2 &= \alpha'x_1 + \beta'x_2 + \gamma'x_3 \\
\lambda x_3 &= \alpha''x_1 + \beta''x_2 + \gamma''x_3,
\end{aligned}}
woraus man durch Elimination von $x_1, x_2, x_3$ die \gesperrt{cubische} Gleichung für $\lambda$ erhält:
\weq{4}{\Lambda = \begin{vmatrix} \alpha - \lambda, & \beta, & \gamma \\ \alpha', & \beta' - \lambda, & \gamma' \\ \alpha'', & \beta'', & \gamma'' - \lambda \end{vmatrix} = 0.}

Es giebt daher im Allgemeinen \gesperrt{drei Punkte}, die ihre eigenen Bilder sind, und diese wollen wir (wie bei den binären Substitutionen) die \gesperrt{Pole der Substitution $A$} nennen. Die drei Verbindungslinien dieser Pole sind dann gleichfalls ihre eigenen Bilder, jedoch so, dass auf jeder dieser drei Geraden nur zwei Punkte liegen, die auf sich selbst abgebildet sind. Die Bilder der übrigen Punkte sind in der Linie verschoben. Diese Linien wollen wir die \gesperrt{Axen der Substitution $A$} nennen.

Es können nun aber besondere Fälle eintreten, die hervorzuheben sind. Es können zwei oder selbst alle drei Pole in einen Punkt zusammenfallen, wie das Beispiel
\[ \begin{pmatrix} \alpha & 0 & 0 \\ \alpha' & \alpha & 0 \\ \alpha'' & \beta'' & \alpha^{-2} \end{pmatrix} \]
zeigt. Hier fallen, wenn $\alpha'$, $\alpha''$, $\beta''$ von Null verschieden sind, zwei, oder wenn $\alpha = 1$ ist, alle drei Pole in einen Punkt zusammen.

Wichtiger noch ist ein anderer besonderer Fall, nämlich der, dass unendlich viele Punkte ihre eigenen Bilder sind. Wenn wir von dem Falle absehen, dass alle Punkte ihre eigenen Bilder sind, der nur bei der Multiplication vorkommt, so müssen die Punkte, die ihre eigenen Bilder sind, wenn ihrer unendlich viele sind, auf einer geraden Linie liegen; denn es kann dieser Fall nur dann eintreten, wenn für eine Wurzel von (4) die drei Gleichungen (3) aus einer von ihnen folgen, oder, was dasselbe ist, wenn mit $\Lambda$ zugleich die sämmtlichen ersten Unterdeterminanten
%CONT>>
